\documentclass[11pt,a4paper]{article}

\usepackage{kotex}
\usepackage{fontspec}
\setmainfont{Times New Roman}
\setmainhangulfont{AppleMyungjo}
\setsanshangulfont{Apple SD Gothic Neo}

\usepackage[a4paper,top=18mm,bottom=18mm,left=20mm,right=20mm]{geometry}
\usepackage{array}
\usepackage{enumitem}
\linespread{1.18}
\setlength{\parindent}{0pt}

\newcommand{\sol}[2]{\vspace{6pt}{\sffamily\bfseries #1} \quad 정답 \fbox{#2}\par\vspace{1pt}}

\begin{document}

\begin{center}
{\fontsize{15}{17}\selectfont\sffamily\bfseries 2026 고1 영어 변형 — 지문 18 정답 및 해설}\\[2pt]
{\sffamily 도서 기부 행사 안내문}
\end{center}
\vspace{6pt}

\begin{center}
\renewcommand{\arraystretch}{1.4}
\begin{tabular}{|c|c|c|c|c|}
\hline
\sffamily 문항 & 1 & 2 & 3 & 4 \\\hline
\sffamily 유형 & \sffamily 목적 & \sffamily 어법 & \sffamily 불일치 & \sffamily 서술형 \\\hline
\sffamily 정답 & ② & ② & ④ & (아래) \\\hline
\end{tabular}
\end{center}
\vspace{8pt}\hrule\vspace{4pt}

\sol{1. 목적}{②}
글쓴이 Trixie Mitchell은 독서 프로그램을 지원하기 위해 중고 도서를 모으는 \textit{book donation drive}에 주민들의 참여(donate)를 요청하고 있다. 따라서 ② ‘to ask residents to donate used books for reading programs’가 정답.
\par\vspace{1pt}\textbf{오답:} ① 운영 시간 연장 공지 아님, ③ 자원봉사 모집 아님, ④ 판매(charity sale)가 아니라 기부, ⑤ 새 도서관 개관 초청 아님.

\sol{2. 어법}{②}
(A) \textit{be pleased to + 동사원형}(\,$\sim$하게 되어 기쁘다\,) → \textbf{to invite}. \quad
(B) ‘$\sim$을 지원하기 위해’라는 목적의 \textit{to부정사} → \textbf{to support}(과거분사 supported는 불가). \quad
(C) 주어 \textit{Any books}가 복수이므로 동사는 \textbf{are}. \par
따라서 (A) to invite / (B) to support / (C) are → 정답 ②.

\sol{3. 일치하지 않는 것}{④}
본문은 ‘\textit{novels, poetry, and \textbf{non-fiction} books}’를 모두 받는다고 했다. ④의 ‘factual books such as biographies(전기 같은 사실 기반 도서)’는 곧 \textit{non-fiction}이므로 수집 대상에 \textbf{포함}된다. 따라서 ‘수집 대상 밖(fall outside)’이라는 ④는 일치하지 않는다.
\par\vspace{1pt}\textbf{확인:} ① responsible for running = director, ② benefit young readers = help young readers discover the joy of reading, ③ worn-out books unlikely welcomed = only \textit{in good condition}, ⑤ drop off when open = during operating hours. 모두 본문과 일치(패러프레이즈).

\vspace{6pt}{\sffamily\bfseries 4. 서술형}\par\vspace{2pt}
(1) \fbox{모범답안} \ \textbf{We gladly accept any books in good condition.}
\par\vspace{1pt}\hspace{1.5em}— 수동태 \textit{are gladly accepted} → 능동태 \textit{accept}로 전환. 목적어 \textit{any books in good condition}, 부사 \textit{gladly}의 위치(동사 앞)에 유의. 베껴 쓸 수 없는 \textbf{태 변환} 문항.
\par\vspace{4pt}
(2) \fbox{모범답안} \ \textbf{이 뜻깊은 행사에 참여해 주신다면 (정말) 감사하겠습니다.}
\par\vspace{1pt}\hspace{1.5em}— \textit{would be grateful if + 주어 + could}\,(가정법)의 공손한 부탁 표현. ‘$\sim$해 주신다면 감사하겠다’로 자연스럽게 해석. \textit{meaningful event} = 뜻깊은/의미 있는 행사.

\end{document}
